\chapter{Setting and main theorem}
\label{chap:setup}

All function spaces are complex valued. Inner products are linear in the
first argument. Nothing in this chapter uses Sobolev spaces: the inverse
Laplacian and the transport operator are defined through Fourier coefficients,
or equivalently distributionally.

\begin{definition}[Torus and Lebesgue spaces]
\label{def:torus}
\lean{NoEmbeddedEigenvalues.L2, NoEmbeddedEigenvalues.char,
  NoEmbeddedEigenvalues.charLp, NoEmbeddedEigenvalues.fourierCoeff,
  NoEmbeddedEigenvalues.MeanZero}
\leanok
Let $\T^2=(\R/2\pi\Z)^2$ with Lebesgue measure $dx\,dy$. Write $L^2=L^2(\T^2;\C)$,
\[
  L^2_0=\Big\{g\in L^2:\int_{\T^2}g\,dx\,dy=0\Big\},
\]
and $\hat g_k=(2\pi)^{-2}\int_{\T^2}g(x)e^{-ik\cdot x}\,dx$, $k\in\Z^2$, for the
Fourier coefficients, so that $g\mapsto(\hat g_k)_k$ is (up to the factor
$2\pi$) a unitary map $L^2\to\ell^2(\Z^2)$ and $g\in L^2_0$ iff $\hat g_0=0$.
\end{definition}

\begin{definition}[Inverse Laplacian]
\label{def:inverse-laplacian}
\uses{def:torus}
\lean{NoEmbeddedEigenvalues.invLap}
\leanok
For $\omega\in L^2_0$ let $\Delta^{-1}\omega\in L^2_0$ be the function with
Fourier coefficients $-\hat\omega_k/|k|^2$ for $k\ne0$ and $0$ for $k=0$. For
$\omega\in L^2_0$ we always write $\psi=\Delta^{-1}\omega$ (the stream function
of the perturbation). Its derivatives up to order two, defined by the Fourier
multipliers $ik_j$ and $-k_jk_l$, are again in $L^2$, and
$\Delta\psi=\omega$.
\end{definition}

\begin{definition}[The dipole and the transport operator]
\label{def:dipole}
\uses{def:torus}
\lean{NoEmbeddedEigenvalues.L0char, NoEmbeddedEigenvalues.HasL0}
\leanok
The steady state is the symmetric dipole with stream function
$\Psi(x,y)=\cos x+\cos y$, vorticity $\Omega=-\Psi$ and velocity
$\nabla^\perp\Psi=(\sin y,-\sin x)$ (manuscript \texttt{(Int1.1)}). Define the
first order differential operators
\[
  L_0=\sin y\,\partial_x-\sin x\,\partial_y,\qquad H=-iL_0,
\]
acting on distributions on $\T^2$. For $g\in L^2$, $L_0g$ is the periodic
distribution with Fourier coefficients ($e_1=(1,0)$, $e_2=(0,1)$)
\[
  \widehat{(L_0g)}_k=\tfrac12\big[k_1(\hat g_{k-e_2}-\hat g_{k+e_2})
                         -k_2(\hat g_{k-e_1}-\hat g_{k+e_1})\big],
\]
and ``$L_0g\in L^2$'' means that this sequence is square summable.
Equivalently, $L_0g=h\in L^2$ if and only if
$\ip{h}{e_k}=-\ip{g}{L_0e_k}$ for all characters
$e_k(x)=e^{ik\cdot x}$, where
$L_0e_k=i(k_1\sin y-k_2\sin x)e_k$; this is the weak definition tested against
characters $\bar e_k$ (as $L_0$ is a real divergence-free field,
$\int(L_0g)\varphi=-\int g\,L_0\varphi$), and expanding $\sin$ into
characters gives the formula above.
\end{definition}

\begin{definition}[Operator domain]
\label{def:domain}
\uses{def:dipole, def:torus}
\lean{NoEmbeddedEigenvalues.Dom}
\leanok
$\Dom=\{g\in L^2_0: L_0g\in L^2\}$ (manuscript \texttt{(Intm3)}).
\end{definition}

\begin{lemma}[The stream function is in the domain]
\label{lem:L0-stream}
\uses{def:inverse-laplacian, def:dipole}
\lean{NoEmbeddedEigenvalues.exists_hasL0_invLap}
\leanok
For $\omega\in L^2_0$, $L_0\Delta^{-1}\omega\in L^2$.
\end{lemma}

\begin{proof}
By the Fourier formula of Definition~\ref{def:dipole},
$|\widehat{(L_0\psi)}_k|\le\tfrac12|k|\sum_{\pm,j}|\hat\psi_{k\pm e_j}|$, and
$|k|\,|\hat\psi_m|\le(|m|+1)|\hat\omega_m|/|m|^2\le2|\hat\omega_m|$ for
$m=k\pm e_j\ne0$. Hence the coefficients are bounded by a finite sum of
shifted copies of $|\hat\omega|$, which is square summable.
\end{proof}

\begin{definition}[Linearized operator]
\label{def:linearized-operator}
\uses{def:domain, def:inverse-laplacian, lem:L0-stream}
\lean{NoEmbeddedEigenvalues.HasL}
\leanok
Let $A=I+\Delta^{-1}$ on $L^2_0$ and $L=L_0A$ with domain $\Dom$; by
Lemma~\ref{lem:L0-stream}, $L\omega\in L^2$ for $\omega\in\Dom$.
\end{definition}

\begin{theorem}[Main theorem]
\label{thm:main}
\uses{def:linearized-operator}
\lean{NoEmbeddedEigenvalues.main}
\leanok
For every $\lambda\in\R\setminus\{0\}$ and $\omega\in\Dom$, if
$L\omega=i\lambda\omega$ then $\omega=0$. That is, $\ker(L-i\lambda)=\{0\}$
(manuscript \texttt{(Intm6)}).
\end{theorem}

\begin{proof}
\uses{lem:reduction, lem:sector-reduction, lem:sector-swap, thm:ee-exclusion, thm:dd-exclusion, thm:mixed-exclusion}
By Lemma~\ref{lem:reduction} it suffices to show that the eigenvalue
problem~\ref{def:eigen-problem} has only the trivial solution. By
Lemma~\ref{lem:sector-reduction} each of the four sector projections of a
solution is again a solution, so it suffices to treat each sector. The
even--even sector is Theorem~\ref{thm:ee-exclusion}, the odd--odd sector is
Theorem~\ref{thm:dd-exclusion}, and the mixed sector $(-,+)$ is
Theorem~\ref{thm:mixed-exclusion}. The other mixed sector $(+,-)$ is mapped to
$(-,+)$ by the rotation $R$ (Lemma~\ref{lem:sector-swap}).
\end{proof}

\begin{definition}[The eigenvalue problem]
\label{def:eigen-problem}
\uses{def:dipole, def:inverse-laplacian}
\lean{NoEmbeddedEigenvalues.IsEigenSolution}
\leanok
A \emph{solution of the eigenvalue problem with parameter
$\lambda\in\R\setminus\{0\}$} is a function $\omega\in L^2_0$ such that, with
$\psi=\Delta^{-1}\omega$,
\[
  H(\omega+\psi)=\lambda\omega
\]
in the sense of distributions (manuscript \texttt{(sf:eigen)}). We say that
$\omega$ is an \emph{eigenfunction}.
\end{definition}

\begin{lemma}[Reduction to the real-parameter equation]
\label{lem:reduction}
\uses{def:eigen-problem, def:linearized-operator}
\lean{NoEmbeddedEigenvalues.hasL_iff_isEigenSolution}
\leanok
Let $\lambda\in\R\setminus\{0\}$. A function $\omega\in\Dom$ satisfies
$L\omega=i\lambda\omega$ if and only if $\omega\in L^2_0$ is a solution of the
eigenvalue problem~\ref{def:eigen-problem} with parameter $\lambda$. Moreover
every such solution lies in $\Dom$.
\end{lemma}

\begin{proof}
\uses{lem:L0-stream}
$L\omega=L_0(\omega+\psi)$ and $H=-iL_0$, so $L\omega=i\lambda\omega$ is
equivalent to $H(\omega+\psi)=\lambda\omega$. If the latter holds in
distributions then $L_0\omega=i\lambda\omega-L_0\psi\in L^2$ by
Lemma~\ref{lem:L0-stream}, so $\omega\in\Dom$.
\end{proof}

\begin{remark}
The manuscript also shows that a nonzero eigenfunction of $L$ can only have
purely imaginary eigenvalue (using $A\ge0$), and explains why $\lambda=0$ must
be excluded: $L$ annihilates $\ker L_0\cap\Dom$ and the first Laplace modes
$\cos x,\sin x,\cos y,\sin y$, which span $\ker A$. Neither fact is needed for
Theorem~\ref{thm:main} as stated.
\end{remark}

\begin{remark}[Lean model]
Decided 2026-10-13, as a first attempt. $L^2$ is Mathlib's \texttt{Lp} for
Lebesgue measure on $\T^2=\texttt{Real.Angle}\times\texttt{Real.Angle}$
(\texttt{NoEmbeddedEigenvalues/Setup.lean}), not $\ell^2(\Z^2)$, because the
later chapters argue in physical space. The operators are defined through
the characters $e_k$:
(i) $L_0$ is a relation, \texttt{HasL0 g h} (``$L_0g=h$''), given by the weak
definition tested against characters (Definition~\ref{def:dipole}); its
equivalence with the Fourier formula is
\texttt{hasL0\_iff\_fourierCoeff}. That smooth test functions give the same
notion is a density argument, deferred until a later chapter needs it;
(ii) $\Delta^{-1}\omega$ is the Fourier series
$\sum_{k\ne0}(-\hat\omega_k/|k|^2)e_k$ in $L^2$, defined on all of $L^2$
(it ignores $\hat\omega_0$);
(iii) the main theorem is stated for $L$ (\texttt{HasL}), and
Lemma~\ref{lem:reduction} connects it to Definition~\ref{def:eigen-problem}.
Mathlib's inner product is conjugate linear in the \emph{first} argument, the
opposite of the convention above.
The remaining infrastructure is Parseval for $\T^2$ with period $2\pi$: Mathlib
has it only for the unit torus (\texttt{mFourierBasis}), so it must be
transported along the scaling map.
\end{remark}
